\chapter{Introduction}\label{intro}

Understanding the history and evolution of the Antarctic Ice Sheet (AIS) requires knowledge of Antarctic bed topography, since it is one of the most crucial boundary 
conditions constraining AIS dynamics. Unfortunately, the direct observation of subglacial landscapes is non-trivial since these are buried under kilometres of ice.
Currently, our ability to observe subglacial terrain relies mostly on indirect methods such as RES and seismic sounding, which return 1D profiles of the bedrock; gravity and magnetic surveys, which constrain composition and basin geometry in 2D; and swath radar, which returns 2D bedrock topography. Even though swath radar represents a major improvement over conventional radar sounding, the typical survey geometries (with line spacings $\geq5$ km and widths of only a couple of kilometres) still yield incomplete coverage and are subject to data degradation by interpolation and down-sampling~\cite{Holschuh_2020}. 

Contemporary maps of the Antarctic topography fall into the following categories: 
\textbf{Interpolated grids}, such as a continent-wide product Bedmap3, which is posted at 500~m, but this is not its resolution, the along-track sampling is metres to kilometres while across-track sampling is hundreds of metres to hundreds of kilometres, so 93\% of its ice-thickness cells are interpolated values~\cite{Pritchard_2025}, so texture at these scales is inferred rather than measured, binding to the recovered bedrock texture by the line spacing
rather than the grid posting. Alternatively, \textbf{inversion models} do not rely on bed measurements but instead infer bed topography indirectly, using ice surface observations and inverse modelling to exploit the fact that basal variability is transmitted by the ice flow to a glacier surface as a function of wavelength~\cite{Gudmundsson_2003}. Inverse methods such as IFPA applied across Antarctica~\cite{Ockenden_2026} are limited in their spatial resolution, since ice flow over bedrock features shorter than the local ice thickness induces negligible perturbations at the ice surface. That floor is a best caseand not the delivered resolution. The IFPA product is presented as a mesoscale map covering 2 to 30~km, with its derived texture metrics: hill count and Fourier fractal dimension computed over wavelengths greater than 5~km and reported on 50~km $\times$ 50~km cells~\cite{Ockenden_2026}. Lastly, there are \textbf{hybrid models}, like BedMachine Antarctica v4~\cite{Morlighem_2026}, are suceptible to both disadvantages described above. Among its methodological strategies to improve on v3, BedMachine Antarctica v4 interpolates the bed by mass conservation in areas of fast ice surface flow (30~m~a$^{-1}$) and by IFPA in areas below that speed threshold. All these examples leave considerable uncertainty in the fine-scale character of the bed.

>>>
It is still uncommon for methods to systematically quantify the relationship between bed topography and its morphometric landscape statistics. Studies that measure the statistics report them as continuous fields 





%JR: I think it is worth noting the value these abilities bring to Antarctic survey design and field planning.